\section{Building Blocks}
\label{sec:method}

\textsc{VisInject} is the composition of two existing papers plus a small evaluation contribution of our own. Figure~\ref{fig:pipeline} shows the resulting three-stage pipeline.

\begin{figure}[t]
  \centering
  \resizebox{\textwidth}{!}{\input{figures/pipeline}}
  \caption{The three-stage \textsc{VisInject} pipeline. Stage~1 trains a single universal adversarial image against $N$ white-box \vlm{}s; Stage~2 transports the resulting signal onto an arbitrary clean photo using a pretrained \textsc{AnyAttack} decoder under an \Linf{} budget; Stage~3 evaluates each (clean, adversarial) response pair along two independent axes.}
  \label{fig:pipeline}
\end{figure}

\subsection{Stage 1 from Rahmatullaev et al.\ \citeyearpar{rahmatullaev2025universal} --- a universal adversarial image}

The first paper trains a single image that, when fed to a target \vlm{} together with \emph{any} of $60$ benign prompts, nudges the \vlm{} toward emitting a fixed target phrase. Concretely, it starts from a grey image, reparameterises pixels as $x = 0.5 + \gamma \tanh(z_1)$ to keep them in $[0, 1]$, and minimises the sum of token-level cross-entropy across all \vlm{}s in a small ensemble and all $60$ prompts. Adam optimiser, $2{,}000$ steps, learning rate $10^{-2}$. The output is a single ``universal adversarial image'' $x_u$ that we use without further modification. We sweep the ensemble size $N \in \{2, 3, 4\}$ as configurations \texttt{2m}, \texttt{3m}, \texttt{4m}.

What we get from this paper is the \emph{signal}: an image with an attack baked in. What is missing is realism --- the universal image looks artificial, so a user uploading it would be flagged.

\subsection{Stage 2 from Zhang et al.\ \citeyearpar{zhang2025anyattack} --- transporting the signal onto any photo}

The second paper trains a CLIP+Decoder pair that maps any image's CLIP feature into an $\eps$-bounded noise tensor. The noise can be added to any clean photo to produce an adversarial photo within an $\Linf$ budget. We use the public weights (\texttt{coco\_bi.pt}) without retraining; we feed our universal image $x_u$ into the encoder and add the resulting noise to each of seven natural test photos:
\begin{equation}
  x_a \;=\; \mathrm{clip}_{[0,1]}\!\Bigl(x_c + \mathrm{clip}_{[-\eps,\eps]}\bigl(\delta(x_u)\bigr)\Bigr), \qquad \eps = 16/255.
\end{equation}
What this paper gives us is the \emph{carrier}: a way to transport an adversarial signal onto a real-looking photo without retraining anything. What is missing is targetedness --- on its own, AnyAttack is good at making the answer drift, but it has no built-in semantic target.

\subsection{Composition --- why it works}

Stage~1 alone makes a strong but obviously-fake adversarial image. Stage~2 alone makes any photo subtly adversarial, but with no specific output target. Composed: Stage~1's universal image already encodes the target phrase in CLIP feature space; passing it through Stage~2's CLIP+Decoder transports that target into noise that can be added to a real photo. The resulting adversarial photo looks like the original ($\Linf = 16/255$, PSNR $\approx 25.2$~dB) and inherits the target signal --- with one important caveat that we measure in \S\ref{sec:results}: the decoder partially erases payload specifics.

\subsection{Stage 3 --- dual-dimension evaluation (our contribution)}

Earlier evaluations \citep{rahmatullaev2025universal} ask a judge model whether the adversarial response ``deviates from'' the clean response. We found this conflates two different things: \textbf{Output Affected} (did the answer change at all?) and \textbf{Target Injected} (did the attacker's chosen phrase actually appear?). On Qwen2.5-VL-3B, our v1 LLM-as-judge reported a $50.5\%$ ``injection rate''; our v2 dual-dim evaluation reports $100\%$ Output Affected but only $0.41\%$ Target Injected on the same data. The difference is not noise --- it is the difference between ``the model lost its train of thought'' and ``the attacker's payload was delivered''.

We replace the judge with two simple programmatic checks: a longest-common-subsequence drift score for Output Affected, and a keyword/regex match (with a small set of close variants per target phrase) for Target Injected. The whole evaluation runs on a laptop CPU and costs nothing per pair --- a non-trivial advantage when sweeping $6{,}615$ pairs.
